\chapter{Architecture}

Nowadays, game makers think of game engines as a single application they download and use to create the game they want. This means that they don't want and need to recompile the engine's source code to develop their own game. This leads to fast iteration times and easier workflows. This means you only need to install the Unity or Unreal editors once, and you can make multiple games without compiling them from source. Of course, if needed most engines do give the ability to modify their source code and create studio/game specific versions of them.

To achieve this, the Game specific code needs to be loaded at runtime, from it's compiled form. In our case, this means that we need to load the DLL containing the game specific code. This is possible on every platform (Windows Linux etc.). But as each platform provides a different set of functions to load in dll files, I will be using the dylib \cite{dylib} library, that makes calling the platform specific functions easier. It provides a single function that wraps each platform's corresponding call, and gives back a pointer that represents the platform specific result.

\begin{minted}{cpp}
dylib lib("foo");

auto adder = lib.get_function<double(double, double)>("adder");

double result = adder(6, 6);
\end{minted}

From here we can then request functions based on their names. However, this can only find functions that were marked exported when the DLL was compiled.


\section{Modules}
To make sure the engine is easy to extend and can be used for any game, it is important to have a mechanism which can add new core components to it. This for example would allow us to have a VR module that we can load in only when we are developing a VR game and not have to keep such a large component in games where it is never going to be used.

To achieve this, we need a way to not include certain parts of the code in the compilation. The simplest way to do this would be to create multiple CMake configurations that leave out the not needed files, or use the preprocessor to disable the code in those sections. The problem with these is that they are hard to maintain, and require the creation of game specific compilations, which goes against the idea of a single executable for the whole engine.

The second solution is to put the potentially not needed code into DLLs and load them at runtime if there is a need. To allow for this, I have created a Module system for Shadow engine. This system covers the startup and core loop of the engine. The engine is split into multiple projects inside the repository, they are each responsible for a specific part of the engine (shown in \autoref{fig:dll-diagram}).

\begin{itemize}
    \item shdaow-engine: The core of the engine that contains most of the code
    \item shadow-runtime: Small startup code for plaiyng a finished game
    \item shadow-editor: The startup and editor related code for the engine
\end{itemize}

\begin{figure}
    \includegraphics[width=1\textwidth]{\auxdir/content/003-architecture/project-layout.eps}
    \caption{The high level components of the engine}
    \label{fig:dll-diagram}
\end{figure}


\section{The App startup}\label{subsec:the-app-startup}
The core of the engine is contained in the shadow-engine project, which gets compiled into a DLL.
That library is loaded in by either the shadow-runtime, or shadow-editor executables, depending on if you are playing or making a game. Both of these have a similar simple main function (\autoref{listing:shadow-runtime-main})

\begin{listing}[!ht]
    \inputminted{cpp}{content/003-architecture/main.cpp}
    \caption{The main function of the runtime executable}
    \label{listing:shadow-runtime-main}
\end{listing}

When starting up, the runtime creates an instance of \verb|ShadowEngine::ShadowApplication|
this is the core of the engine, located in the shadow-engine DLL. This class owns the Module Manager, which is responsible for loading in the modules that make up the engine. Every module needs to extend the base \verb|Module| class.

In the \verb|Init| function, the application does all the module loading that is necessary before the main loop starts. For this it needs to load in a few assemblies\footnote{I'm calling the DLLs assemblies in the code as this is the naming C\# uses. They refer to slightly different concepts, but for this purpose can be treated as the same.}. These are the actual assembly of the \verb|shadow-engine| DLL, the assembly the game is in, and the editor assembly if that is to be used. (shown in \autoref{fig:startup-sequence})

After this, each module needs to be registered to the Module manager using a descriptor (\autoref{listing:module-descriptor}). This object describes everything that is needed to create the module, and what other modules it needs to function correctly. The \verb|id| field is a string identifier for the module \footnote{This ID format is compatible with the planned resource system, allowing us to plater treat the modules as a first class resource if needed.}, and the \verb|assembly| field is a similar ID that defines in witch DLL the manager needs to search for this class. The \verb|dependencies| is a list of string IDs of the modules that need to be loaded before this module.

\begin{listing}[!ht]
    \begin{minted}{cpp}
m.AddDescriptors({
	 .id="module:/renderer/vulkan",
	 .name = "Vulkan",
	 .class_name = "VulkanModule",
	 .assembly="assembly:/core",
	 .dependencies={"module:/platform/sdl2"},
 });
    \end{minted}
    \caption{Example of a module descriptor}
    \label{listing:module-descriptor}
\end{listing}

This registration is done by calling a specially named function in each assembly. These functions must have the name \verb|assembly\_entry|, this is used as a known entry point into each DLL. These functions are called by the \verb|LoadModulesFromAssembly| function. But in the future, we plan on moving the module descriptors fully out of the code and treating them as text based assets.

After both all the assemblies are added and the modules from them are registered, the app calls the module manager's \verb|Init| function (shown in \autoref{fig:startup-module-init}). This, based on the dependencies described, uses a Depth First Search to try and create a topological sort. Currently, it is the user's responsibility to ensure that the dependencies are cycle free. When this order is found, the module manager tries to create each of the modules. This is done by searching for a special function in the assembly the module is located in. This function needs to follow these rules:
\begin{itemize}
    \item Must be marked with \verb|extern "C"|
    \item It's name must be in the format \verb|<class name>_entry()| where the \verb|<class name>| section is the same as the classes name without the namespace section.
    \item Needs to give back a \verb|std::shared_ptr| to a freshly created instance of the module.
\end{itemize}

To help create these functions I have made \verb|MODULE_ENTRY| which is a preprocessor macro that creates them based on the class name.

After every module was successfully instantiated, the module manager calls pre-Init and Init on the modules. Between these steps, the DFS sorting is re-run to ensure every module still has the necessary modules to continue loading. For example, a VR module might exit in pre-init if it finds no VR headset to work with.

\begin{figure}
    \includegraphics[width=1\textwidth]{\auxdir/content/003-architecture/startup.eps}
    \caption{The high level components of the engine}
    \label{fig:startup-sequence}
\end{figure}

\begin{figure}
    \includegraphics[width=1\textwidth]{\auxdir/content/003-architecture/startup-init.eps}
    \caption{The high level components of the engine}
    \label{fig:startup-module-init}
\end{figure}

After this initialization sequence is done, the runtime calls the Start method on the shadow application which starts the main update loop.


\section{The update loop}

